% !TeX root = ../main_en.tex
\subsection{Quadcopter Simulation and Modeling}

The quadcopter is a six-degree-of-freedom aerial vehicle whose position and attitude evolve in a coupled manner, and whose actuation originates from a limited number of rotors. This combination makes trajectory tracking a particularly notable problem within unmanned systems, since the vehicle must modify its orientation to achieve the required translation while maintaining stability and respecting actuation limits. For this reason, literature on small UAVs tends to describe the problem using rigid-body models, well-defined reference frames, and cascaded control architectures \cite{beard2012small,mellinger2011minimum,lee2010geometric}.

In a 6DOF rigid-body model, the vehicle state combines position, velocity, attitude, and angular velocity. Translation is achieved by summing gravity, aerodynamic forces, and the force generated by the rotors, transformed from the body reference frame to the local frame. Rotation is described by Euler's equations, where the moments applied by the rotors modify the angular velocity as a function of the inertia matrix. This type of formulation allows studying the interaction between position tracking, tilt, actuator limitations, and stability, although it does not reproduce all the aerodynamic phenomena of flight.

The term \textit{6DOF} indicates that the vehicle has three degrees of freedom in translation and three in rotation. However, it cannot independently impose accelerations on each axis and arbitrary attitude at the same time. In a traditional quadcopter, the rotors generate thrust along the vertical axis of the vehicle. Lateral force appears by tilting the quadcopter and, thereby, orienting a portion of the thrust toward the desired horizontal direction. This is the physical reason why position tracking and attitude control are said to be coupled.

Reference frames are an essential part of this formulation. In aircraft and aerial vehicles, a distinction is made between a reference frame associated with the environment and another associated with the vehicle. The former allows expressing position, velocity, and trajectory references in an absolute manner, while the latter facilitates determining the moments and forces generated by the rotors. The transformation between the two is performed via a rotation matrix associated with attitude. In this context, quaternions are a common alternative to Euler angles as they represent orientation without introducing kinematic singularities at specific attitudes \cite{kuipers1999quaternions}. Euler angles remain useful for visualization and for setting operational limits, but they are not the most robust method for attitude computation and propagation during numerical integration.

Force generation in a quadcopter is modeled by a quadratic relationship between the rotor angular velocity and thrust. Each rotor also contributes to the total moment through its position relative to the center of mass and the reaction torque associated with its rotational direction \cite{beard2012small,pounds2010modelling}. The actuator mixer solves the inverse problem: given a total thrust and three moments, it calculates the individual thrusts that the rotors should produce. When the solution requires negative thrusts or thrusts exceeding the maximum available, command saturations occur. In this case, the controller may have calculated a coherent yet unrealizable command. Therefore, the comparison of controllers is not limited to the position error, but also records when an actuation could not be carried out.

\begin{figure}[H]
  \centering
  \includegraphics[width=1 \textwidth]{FIG-015}
  \caption{Actuation and motion mechanisms in a quadcopter. For each maneuver, the top view of the vehicle in layout (left) is shown with the rotational direction of each rotor and indications of thrust imbalance ($+$, $-$, $=$), and the resulting physical effect in side view (right). Tilting the vehicle ($\phi$, $\theta$) redirects the total thrust vector ($T$), generating the lateral ($F_y$) or longitudinal ($F_x$) force components that enable horizontal translation. Source: own elaboration.}
  \label{fig:actuacion-movimiento}
\end{figure}

Actuation management can be understood from four basic combinations (see Figure~\ref{fig:actuacion-movimiento}). If all four rotors increase or decrease their thrust in a similar manner, the total thrust changes, and therefore the vertical acceleration in a level attitude changes. If thrust increases on one side relative to the opposite side, a roll moment appears, and if the difference occurs between the front and rear parts, a pitch moment appears. Yaw is obtained by leveraging the reaction torques of rotors spinning in opposite directions. In summary, each rotor contributes a force $T_i$, and the assembly produces an output of the type
\[
  \bm{u} = [T\;\tau_x\;\tau_y\;\tau_z]^\mathsf{T},
\]
where $T$ is the collective thrust and $\tau_x$, $\tau_y$, and $\tau_z$ are the moments in the FRD body frame.

The time integration of the model requires discretizing the differential equations. Runge--Kutta methods, and in particular RK4, are a classical option when seeking a favorable trade-off between precision and computational cost \cite{butcher2016numerical}. In digital control environments, it is common to separate the physics integration frequency, the controller update frequency, and the telemetry logging frequency. Between two control cycles, the last calculated command is kept constant using a zero-order hold (ZOH), thus approximating the behavior of a digital system acting on continuous dynamics \cite{astrom1997computer}.

Proven platforms exist for simulating aerial vehicles. MATLAB, Simulink, and the UAV Toolbox provide tools to design, simulate, test, and deploy UAV applications, including scenarios, sensors, autopilots, and test workflows both on a computer and in hardware-in-the-loop configurations \cite{mathworksUAVToolbox}. In the open-source domain, RotorPy offers a Python multirotor simulator oriented toward education and research, featuring 6DOF dynamics, aerodynamics, actuators, sensors, wind, and validation against real data \cite{folk2023rotorpy}.

The present work does not need to reproduce the full breadth of these platforms; instead, complete modularity and traceability are the most important aspects. The previously discussed parts are implemented: 6DOF dynamics, specific reference frames, attitude using quaternions, quadratic control law, mixer with saturation, simplified disturbances, integration, and sufficient telemetry to reproduce and review each run.

\subsection{Programming and Machine Learning Environments}

The choice of an implementation environment conditions the available tools, reproducibility, and ease of inspection of the work. Python has established itself as an environment for scientific programming and simulations, with an open license and an extensive ecosystem of libraries. Its interest in this context does not lie in being the most efficient option for a future onboard flight, but rather in allowing the model to be expressed clearly, generating data, running comparisons, and connecting the simulator with machine learning tools.

PyTorch provides tensors, loss functions, and optimizers that allow training supervised and recurrent models on CPU or GPU \cite{pytorchDocs}. Furthermore, it includes standard implementations of linear layers, GRU, LSTM, MSE loss, and Adam, allowing work at a higher level and avoiding the risk of introducing errors in widely established components.

Reproducibility also depends on environment management and the files related to experiments. Tools such as \texttt{uv} allow declaring, installing, and running Python environments in a repeatable and simple manner \cite{uvDocs}. The use of human-readable configuration files, such as YAML (\textit{YAML Ain't Markup Language}, a readable text format to describe structured data using keys, lists, and values), makes it easy for scenarios, seeds, trajectories, controllers, and termination conditions to be reviewed without modifying code in a visual and simple way. Compared to JSON, which is widespread but more rigid for reading and writing by a researcher, and TOML, which is clear for relatively flat configurations, YAML offers a convenient syntax for describing hierarchical structures with lists, dictionaries, and nested blocks \cite{yamlSpec,jsonRFC,tomlSpec}. Finally, Git and GitHub provide common support for change history, code distribution, and traceability of experiments \cite{githubDocs}. Thanks to this toolset, auditing and reproducing experiments by anyone is facilitated.

\subsection{Classic Cascaded Control and Gain Scheduling}

Classic quadcopter control is typically organized in cascade. The outer loop works on position and velocity, and calculates a desired force or acceleration in the ENU world reference frame. The inner loop transforms this translational intent into a desired attitude and stabilizes the vehicle using moments in the FRD body reference frame. This separation is also observed in speed: translational dynamics are slower and require fewer corrections than angular dynamics \cite{beard2012small,mellinger2011minimum,lee2010geometric}.

In this architecture, the outer loop does not directly decide the rotor speeds; rather, its output is a desired force in the local frame. If the vehicle is below the reference, the force will have a larger vertical component; and if it needs to move horizontally, it will have a component in this direction. The inner loop takes the force, calculates what vehicle orientation would allow generating it with the available thrust, and modifies the attitude toward that orientation. The reference yaw completes the orientation, since the thrust direction only fixes the direction of the vertical axis.

PD and PID controllers remain the benchmark due to their simplicity, interpretability, and low computational cost. A proportional term corrects the position or attitude error, a derivative term dampens the response based on velocity or angular velocity, and an integral term can reduce steady-state errors. However, the integral action can accumulate command during saturations, thus requiring anti-windup protection mechanisms. In this case, a PD structure combined with trajectory acceleration feedforward and gravity force compensation results in a more feasible implementation.

Gain tuning can be approached in several ways. Manual tuning provides intuition but hinders traceability if all iterations are not preserved. Classic Ziegler--Nichols rules offer systematic procedures for certain classes of systems, although they do not directly represent a 6DOF quadcopter with saturations and spatial trajectories \cite{ziegler1942optimum}. It is also possible to use controllers on models linearized around operating points, or to search for gains using numerical optimization evaluated in closed-loop simulation. This latter approach is of special interest when the objective is to compare various candidates under the same trajectories, limits, and metrics.

When a single set of gains does not yield a satisfactory result across the entire operating range, a classic response is gain scheduling. This technique selects or interpolates control parameters according to certain variables describing the operating point, so that different flight regions can use different tunings \cite{leith2000survey}. Conceptually, a set of specialized controllers by trajectory family is related to this idea, as there are several experts and the most appropriate one is applied for each condition. However, in this work, the neural implementation does not perform an explicit gain interpolation according to fixed criteria; instead, it learns to predict the force of the outer loop from the same data used by the PD.

\subsection{Imitation Learning and Hybrid Control}

Imitation learning aims to train a system using data generated by an expert system. In its most direct form, this behavior cloning poses a supervised problem: the system's observations are received, and it learns to reproduce its actions. This concept has antecedents in autonomous driving, such as ALVINN, where a network learned to produce steering commands from human examples \cite{pomerleau1989alvinn}. In robotics and control, imitation avoids the design of a complete reward function and allows behavior transfer from an already available controller.

However, behavior cloning has a significant limitation: the data comes from states visited by the expert, but during the execution of the learned policy, small errors can accumulate, leading to states not represented in the training data. This discrepancy between the training distribution and the one induced by the learned system is known as distribution shift (covariate shift). Methods such as DAgger \cite{ross2011dagger} reduce this problem by aggregating data from states visited by the neural system and querying the original expert system again. Although this work does not implement DAgger, it is important to remember that a low supervised loss during training does not guarantee good behavior in closed loop.

In physical control applications, a machine learning policy can replace the entire controller or only a part of the architecture. Hybrid approaches preserve classic elements when these contribute stability, limits, or interpretability. In this case, a network learns the behavior of the outer loop, while the inner loop, mixer, and protections remain classic. However, this introduces the risk of erroneously attributing effects to the network, meaning the contributions of each element must be identified and separated.

\subsection{Neural Networks for Control}

The neural networks used in this work are formulated as supervised regression models. The input is a vector or sequence of vectors constructed from the observation and the trajectory reference, while the output is a three-dimensional force in the ENU world reference frame. During training, the parameters are adjusted to reduce the difference between the predicted force and the force commanded by the specialized PD. This formulation requires representing the data as tensors, grouping samples into batches, and optimizing a differentiable loss function via backpropagation \cite{rumelhart1986learning,pytorchDocs}.

The training of a network is not determined solely by the architecture and the data. It also depends on hyperparameters---decisions set before the tuning of weights and biases that govern how the optimization is performed. Among these are the learning rate, batch size, maximum number of epochs, and stopping criteria. Unlike the trainable parameters of the network, these values are not learned but are selected as part of the design.

A batch is the set of samples processed together before calculating a loss and updating the model parameters. Its size is usually denoted by $B$; larger values provide a more stable estimate of the gradient, while smaller values introduce more variation between updates and require more iterations. An epoch corresponds to a complete pass through the training samples. Early stopping halts training when the validation loss stops improving, and patience indicates how many consecutive epochs without improvement are tolerated before interrupting the process.

A neural network does not store explicit rules of the \textit{if--then} type; instead, it learns weights and biases that transform inputs into outputs. In supervised regression, each sample provides an input--target pair, and training modifies the parameters to minimize a loss. If $\theta$ groups all the parameters of the network and $\mathcal{L}(\theta)$ measures the average error, gradient descent updates the parameters in the direction opposite to the gradient:
\[
  \theta_{n+1}=\theta_n-\eta\nabla_\theta\mathcal{L}(\theta_n),
\]
where $\eta$ is the learning rate. A value that is too large can cause the training to oscillate or diverge; one that is too small can make it excessively slow. Optimizers like Adam adapt the effective update step size using accumulated gradient information, but they still depend on a base learning rate \cite{kingma2015adam}.

\subsubsection{Feedforward Networks}

A feedforward network computes the output through a succession of layers applied in a single direction, from the input to the output prediction. In a linear layer, the output is obtained by multiplying the input by a weight matrix and adding a bias vector. If only linear transformations were chained, the model would remain equivalent to a global linear transformation. This is why non-linear activation functions are introduced between the layers. The Rectified Linear Unit (ReLU) is a common activation due to its simplicity and because it facilitates training compared to the saturations inherent in other functions \cite{nair2010relu}.

For a general layer, the operation can be written as
\[
  \bm{z}^{(\ell)}=\bm{W}^{(\ell)}\bm{a}^{(\ell-1)}+\bm{b}^{(\ell)},\qquad
  \bm{a}^{(\ell)}=\phi(\bm{z}^{(\ell)}),
\]
where $\bm{W}^{(\ell)}$ and $\bm{b}^{(\ell)}$ are weights and biases, and $\phi$ is a non-linear activation. Depth is given by the number of layers, while width is associated with the number of units per layer.

The multilayer perceptron, or MLP, combines several linear layers with non-linear activations. Theoretically, feedforward networks with enough units can approximate a wide range of continuous functions \cite{cybenko1989approximation}. Thus, the MLP is used as the memoryless neural baseline, so that two identical inputs produce the same force regardless of what occurred in the previous step. This feature makes it simpler and computationally cheaper than recurrent architectures, but it also limits its ability to use temporal information when lags or transient states exist.

\subsubsection{Recurrent Networks}

Recurrent neural networks (RNNs) extend the feedforward formulation through a hidden state that is updated over time. At each step, the model processes the current input and a compact representation of the prior history. Thus, the output depends not only on the error at that instant, but also on how that state was reached.

A basic RNN can be summarized by
\[
  \bm{h}_k=\phi(\bm{W}_x\bm{x}_k+\bm{W}_h\bm{h}_{k-1}+\bm{b}_h),\qquad
  \bm{y}_k=\bm{W}_y\bm{h}_k+\bm{b}_y.
\]
The vector $\bm{x}_k$ is the input at the given step $k$, $\bm{h}_k$ is the hidden state summarizing the processed history up to that step, and $\bm{y}_k$ is the output. During training, the network is unrolled in time, and the gradient is backpropagated through the temporal sequence. This procedure, known as backpropagation through time (BPTT), allows tuning the same parameters from the errors observed in the temporal sequence.

However, the basic RNN presents difficulties in learning long temporal dependencies. During backpropagation, gradients can vanish or explode, hindering the preservation of relevant information over a large number of steps. This is why modern recurrent architectures feature internal gates to regulate which information is kept, updated, or discarded \cite{bengio1994longterm}. Figure~\ref{fig:comparativa-arquitecturas-control} illustrates the fundamental difference between the static and temporal approaches of these two types of architectures.

\paragraph*{Long short-term memory (LSTM).}
The LSTM architecture introduces a cell state and several gates controlling the input, forgetting, and output of information. Its goal is to facilitate the network's retention of relevant information over more time steps than a simple RNN, mitigating the problem of vanishing gradients \cite{hochreiter1997lstm}. In the context of trajectory tracking, this memory is beneficial if the recent history contains information about delays or error evolution. However, it also causes an increase in network parameters and can introduce sensitivity to temporal patterns that are poorly represented in the training dataset.

\begin{figure}[H]
  \centering
  \includegraphics[width=\textwidth]{FIG-014}
  \caption{Conceptual comparison of neural architectures for control. In (A), the purely static nature of the MLP is shown. In (B), the unrolled flow of the basic RNN and its vulnerability to gradient vanishing are detailed. In (C) and (D), the selective gated structure of the LSTM cell and the GRU cell are contrasted. Source: own elaboration.}
  \label{fig:comparativa-arquitecturas-control}
\end{figure}

The LSTM separates two internal quantities. The hidden state $\bm{h}_k$ is the representation output from the cell, used to calculate the prediction or serve as input for the next step. The cell state $\bm{c}_k$ acts as a more direct internal memory, designed to carry information along the sequence. The three main gates regulate this memory: the forget gate decides what portion of $\bm{c}_{k-1}$ is kept, the input gate decides what new information is incorporated, and the output gate decides what portion of the updated memory is exposed as $\bm{h}_k$. A common formulation of the LSTM is
\begin{align}
  \bm{f}_k &= \sigma(\bm{W}_f[\bm{h}_{k-1},\bm{x}_k]+\bm{b}_f),\\
  \bm{i}_k &= \sigma(\bm{W}_i[\bm{h}_{k-1},\bm{x}_k]+\bm{b}_i),\\
  \tilde{\bm{c}}_k &= \tanh(\bm{W}_c[\bm{h}_{k-1},\bm{x}_k]+\bm{b}_c),\\
  \bm{c}_k &= \bm{f}_k\odot\bm{c}_{k-1}+\bm{i}_k\odot\tilde{\bm{c}}_k,\\
  \bm{o}_k &= \sigma(\bm{W}_o[\bm{h}_{k-1},\bm{x}_k]+\bm{b}_o),\\
  \bm{h}_k &= \bm{o}_k\odot\tanh(\bm{c}_k).
\end{align}
Here, $\sigma$ is the sigmoid activation, $\tanh$ is the hyperbolic tangent activation, $[\bm{h}_{k-1},\bm{x}_k]$ denotes concatenation, and $\odot$ is the Hadamard (element-wise) product.

\paragraph*{Gated recurrent unit (GRU).}
The GRU also uses gates, but with a more compact structure than the LSTM. Instead of explicitly separating a cell state and a hidden state, it combines update and reset mechanisms to decide how much prior information to retain and how to incorporate the current input \cite{cho2014learning}. This positions it as a recurrent network alternative with less complexity than the LSTM while maintaining sequence-processing capabilities. Thus, comparing GRU and LSTM under the same conditions will allow assessing whether a more complex memory provides an improvement for outer-loop control.

In a GRU, the update gate $\bm{z}_k$ regulates how much of the previous state is kept versus how much is replaced by a new proposal. The reset gate $\bm{r}_k$ controls how much the previous state influences the candidate calculation $\tilde{\bm{h}}_k$. A common formulation is
\begin{align}
  \bm{z}_k &= \sigma(\bm{W}_z[\bm{h}_{k-1},\bm{x}_k]+\bm{b}_z),\\
  \bm{r}_k &= \sigma(\bm{W}_r[\bm{h}_{k-1},\bm{x}_k]+\bm{b}_r),\\
  \tilde{\bm{h}}_k &=
    \tanh(\bm{W}_h[\bm{r}_k\odot\bm{h}_{k-1},\bm{x}_k]+\bm{b}_h),\\
  \bm{h}_k &= (1-\bm{z}_k)\odot\bm{h}_{k-1}
             +\bm{z}_k\odot\tilde{\bm{h}}_k.
\end{align}
By not maintaining a separate cell state, the GRU typically requires fewer parameters than an equivalent LSTM. However, there may be cases where this lower complexity is sufficient to capture the temporal dependency and even reduce sensitivity to sequences with low representation.

In inference, there are two practical ways to use a recurrent network. In a \emph{stateful} execution, the hidden state computed in one step is preserved and supplied as the initial state of the next step. In a window-based execution, by contrast, for each prediction a sequence of fixed length is reconstructed, fully processed, and the output is provided. The first option avoids repeatedly recalculating past information but causes a dependence on the network's internal propagation state throughout the episode. The second is computationally more expensive but maintains a local impact restricted to the window. This distinction will be relevant later, as the implementation carried out uses fixed-length temporal windows instead of a fully stateful recurrent inference.

The training of these architectures retains elements of supervised learning. Normalization calculated solely using the training split prevents variables of different scales from dominating the loss function and reduces information leakage from validation or testing. The mean squared error (MSE) is a useful function in this case, penalizing large deviations between prediction and ground truth with greater intensity. Validation is also used to select a checkpoint without tuning weights on the test set, and early stopping is applied to halt training if no improvement is observed after a set number of epochs.

\subsection{Evaluation, Generalization and Metrics}

In dynamic systems of this type, evaluation must be conducted both in a supervised manner and on the behavior of the system in closed loop. It can happen that a network approximates the specialized controller's force very well on test set samples and yet generates a worse or even divergent trajectory when its predictions are fed back into the dynamics. This discrepancy stems from the fact that closed-loop control modifies the future states seen by the network itself, turning deviations in forces into different input data. For this reason, an evaluation of imitation learning must distinguish between supervised fidelity, tracking error, and the utilization of system limitations.

The separation between training, validation, and testing is necessary to estimate performance on data not used for parameter tuning or checkpoint selection. Furthermore, when studying transfer, one must also distinguish between in-distribution (ID) and out-of-distribution (OOD) conditions. Thus, the former verify interpolation within families and disturbances already represented in training, while the latter evaluate whether a reasonable behavior is preserved under recombinations or new geometries.

Tracking metrics must capture complementary aspects of the error. In this way, RMSE penalizes large errors more heavily, MAE describes the average error, and maximum error detects localized extremes. These quantities must be interpreted alongside saturations, control effort, and simulation termination causes, as a low average error can be achieved with a continuous demand for unrealizable actions.

\subsection{Positioning of This Work}

This work positions itself at an intermediate point between classic control and purely learned control. On one hand, cascaded control allows the separation between translation and attitude, observability of its functioning, force, moment, and saturations, but faces a limitation due to the tuning process depending on the trajectory family, transient motions, and disturbances considered. On the other hand, neural networks are a flexible way to approximate relationships from data, but require a distinction between supervised fidelity and closed-loop behavior.

This work, therefore, employs a hybrid scheme where the classic inner loop, the mixer, and the actuation protections are maintained, and only the outer loop is replaced by an imitation-trained neural network. The network generates a desired force in the ENU world frame, thus preserving the separation between the stochastic component and the deterministic elements. In this context, a simulator with traceability is constructed, specialized PD controllers are tuned, and MLP, GRU, and LSTM architectures are trained to condense the behavior of classic controllers. The comparison supplements the supervised loss with tracking metrics, terminations, and closed-loop protection activation.
